\documentclass[11pt,a4paper]{article}
\usepackage[T2A]{fontenc}
\usepackage[utf8]{inputenc}
\usepackage[russian,english]{babel}
\usepackage{amsmath,amssymb,amsthm,mathtools}
\usepackage{setspace}
\usepackage{enumitem}
\usepackage[hidelinks]{hyperref}
\setstretch{1.3}
\frenchspacing
\hyphenpenalty=10000
\renewcommand{\leq}{\leqslant}
\renewcommand{\le}{\leqslant}
\renewcommand{\geq}{\geqslant}
\renewcommand{\ge}{\geqslant}
\newcommand{\head}[1]{\par\vspace{15pt}\noindent\textbf{\underline{#1}}\quad\ignorespaces}
\newcommand{\src}[1]{\unskip\nobreak\hfil\penalty50\hskip1em\hbox{}\nobreak\hfill(#1){\parfillskip=0pt\par}}
\begin{document}
\begin{center}
{\huge\bfseries Methods of proof -- 4\par}\vspace{1.1em}
{\Large Invariants, monovariants and double counting\par}\vspace{1.1em}
{\Large Homework\par}\vspace{1.1em}
{\foreignlanguage{russian}{13 сентября 2026 года}\par}
\end{center}\vspace{0.0pt}

\noindent\textbf{\underline{Theoretical minimum.}}\quad An invariant proves impossibility, a
monovariant proves termination, and double counting proves identities and inequalities. To find an
invariant, look at what the allowed operation does to the simplest functionals: sum, sum of
squares, parity, residue, sign of a permutation, rank, determinant; the right invariant is usually
the one that the operation multiplies by a constant or leaves untouched. To find a monovariant,
look for a bounded integer quantity that strictly decreases, or for a continuous quantity with a
sign-definite derivative. Double counting is the habit of computing one number in two ways; when
the count is over incidences it produces an inequality, and when it is over a set of pairs it
produces a bijective proof. A generating function is the same bookkeeping written algebraically.

\medskip
\noindent\rule{\textwidth}{0.4pt}

\begin{enumerate}[itemsep=0.4em,topsep=7.4pt]
\item Numbers $1,2,\dots,n$ are written on a board; one may erase two numbers and write their
difference. Determine for which $n$ the last remaining number can be zero, and find the invariant
responsible.

\item Prove that a $10{\times}10$ board cannot be tiled by $1{\times}4$ pieces, by a colouring
invariant, and determine all $m\times n$ boards tileable by $1\times k$ pieces.

\item In a group of people every two are friends or enemies. Each step, a person with more
enemies than friends changes sides: his friends become his enemies and vice versa. Prove that the process terminates, by exhibiting a monovariant, and estimate the
number of steps.

\item Prove that $\sum_k\binom nk^2=\binom{2n}n$ by double counting, and
$\sum_kk\binom nk=n2^{n-1}$ in two ways; then give a bijective proof of Vandermonde's identity.

\item Prove that in any finite set system the number of pairs $(x,A)$ with $x\in A$ counted in two
ways gives Fisher's inequality for a design, and deduce the bound on the number of blocks.

\item Prove that the trace of a nilpotent matrix is zero by using an invariant of similarity, and
deduce that $AB-BA=E$ is impossible over a field of characteristic zero (compare with Linear
algebra -- 12).

\item Prove that the parity of the number of inversions is invariant under a transposition, and
use it to show that the alternating group has index two; then apply the sign invariant to a
puzzle of your choice.

\item A finite set of points in the plane is such that the midpoint of any two of them also
belongs to the set. Prove that the set has one point, and give the monovariant proof and the
extremal proof; compare with the lattice version (Convex -- 10).

\item Prove that a sequence of moves in which a number is replaced by the sum of its digits
terminates, find the invariant modulo $9$, and prove the general statement about the digital root.

\item Prove Euler's formula $V-E+F=2$ for planar graphs by an induction that is really a
monovariant argument, and deduce that $K_5$ is not planar by double counting the incidences
between faces and edges.
\end{enumerate}
\end{document}
